\section{Related Work}

\subsection{Linear-Complexity Attention}

\textbf{Linformer.} Wang et al. (2020) project $K, V$ to fixed low rank with 
learned projection matrices. Galerkin uses a \textbf{fixed DCT basis}, not learned.

\textbf{Performer.} Choromanski et al. (2021) approximate softmax via random 
features. Galerkin computes \textbf{exact} attention in basis space.

\textbf{Nystr mformer.} Xiong et al. (2021) use Nystr m landmarks. Galerkin uses 
a full DCT basis.

\textbf{Linear Transformer.} Katharopoulos et al. (2020) reorder attention as 
associative products. Different mathematical primitive.

\subsection{Basis Projection Methods}

\textbf{FNet.} Lee-Thorp et al. (2021) replace self-attention with an unparameterized 
Fourier transform. Galerkin uses a truncated DCT basis plus a softmax 
attention in basis space.

\textbf{GFNet.} Rao et al. (2021) use global filter operations in Fourier space 
for vision. Similar spirit, different domain.

\textbf{AFNO.} Guibas et al. (2022) use adaptive Fourier neural operators for 
PDEs. Galerkin is a Galerkin discretization of the attention operator.

\subsection{Numerical Analysis of Attention}

\textbf{Neural Operators.} Kovachki et al. (2023) formalize neural operators as 
mappings between function spaces. Attention is a specific case.

\textbf{Galerkin Methods.} Classical theory (Strang, 1973; Brenner \& Scott, 2008) 
for variational discretization of PDEs.

\textbf{Petrov-Galerkin.} Uses different test and trial bases. Galerkin Attention 
uses a \textbf{single} orthonormal basis (Galerkin, not Petrov-Galerkin in the 
strict sense).

\subsection{Position of Our Work}

\begin{table}[H]
\centering
\small
\begin{tabular}{lllll}
\toprule
\textbf{Method} & \textbf{Basis} & \textbf{Learned} & \textbf{Encoder} & \textbf{Decoder} \\
\midrule
Linformer & Linear proj & Yes & Yes & Yes \\
Performer & Random feats & No & Yes & Yes \\
FNet & FFT & No & Yes & No \\
Nystr mformer & Landmarks & Yes & Yes & Yes \\
\textbf{Galerkin} & \textbf{DCT (fixed)} & \textbf{No} & \textbf{Yes} & \textbf{Fallback} \\
\bottomrule
\end{tabular}
\end{table}

\textbf{Key distinction.} Galerkin is the only method that:
\begin{enumerate}
\item Uses a \textbf{fixed orthonormal basis} (DCT).
\item Dispatches encoder vs decoder.
\item Provides an \textbf{exact} solution in basis space.

\end{enumerate}
\subsection{Comparison Table}
